\documentclass[11pt]{article}
\usepackage[margin=1in]{geometry}
\usepackage{amsmath,amssymb,amsthm}
\usepackage{booktabs,longtable,array}
\usepackage[hidelinks]{hyperref}

\title{Möbius-Admissible Completions, Short Correlations, and RH-Scale Excursions:\\A Literature and Research Map}
\author{Víctor Casaña-Giner\\Bio Logic Crop Science S.L., E-46120 Alboraya (Valencia), Spain}
\date{August 2026}

\newcommand{\Om}{\Omega}
\newcommand{\RHC}{\mathrm{RHC}}
\newcommand{\RHF}{\mathrm{RHF}}
\newcommand{\RHFl}{\mathrm{RHFl}}
\newcommand{\BN}{B_N}

\begin{document}
\maketitle

\begin{abstract}
This document is a research map for a deterministic extremal program around the Möbius summatory function. The starting space is the full ternary cube $\{-1,0,1\}^N$. Successive models impose explicit pieces of Möbius information and ask what endpoint sums and short correlations are still possible among all surviving completions. The current program has two connected parts. The first studies contraction of the adversarial endpoint envelope under arithmetic information. The second studies whether a large endpoint excursion forces an early change in the short-correlation vector relative to the true Möbius word. The document records the exact finite models, closed-form baseline results, locality barriers, short-correlation identities, the current exhaustive-computation specification, and the relevant literature. The individual ingredients belong to established areas such as multiplicative functions, discrepancy, random multiplicative functions, sign patterns, and Mertens-function theory. The specific finite completion and short-correlation questions are therefore stated with explicit rule systems and with novelty claims kept separate from the mathematical definitions.
\end{abstract}

\section{The organizing problem}
Let
\[
 \Om_N^{(0)}=\{-1,0,1\}^N,\qquad |\Om_N^{(0)}|=3^N,
\]
and let $\Om_N^{(j)}\subseteq\Om_N^{(0)}$ be the set of words surviving a specified collection $\mathcal I_j$ of rigorous arithmetic constraints.  For $\omega=(\omega_1,\ldots,\omega_N)$ put
\[
 S_\omega(n)=\sum_{k\le n}\omega_k,
 \qquad
 B_N^{(j)}=\max_{\omega\in\Om_N^{(j)}}|S_\omega(N)|.
\]
The quantity
\[
 E_N^{(j)}=\frac{\log B_N^{(j)}}{\log N}-\frac12
\]
was used in the first finite tables, but it is not a reliable finite-scale estimate of the growth exponent. If
\[
 B_N=cN^\theta,
\]
then
\[
 E_N=\theta-\frac12+\frac{\log c}{\log N},
\]
so a small constant $c$ can make a linear envelope look close to square-root scale at small $N$. The main diagnostics are therefore
\[
 R_N^{(j)}=\frac{B_N^{(j)}}{N}
\]
and the local logarithmic slope
\[
 \theta_N^{(a,j)}=
 \frac{\log B_{aN}^{(j)}-\log B_N^{(j)}}{\log a},
\]
with $E_N^{(j)}$ retained only as a historical finite statistic. The central deterministic question is not whether a large excursion is improbable, but whether any admissible completion producing it remains.

A useful filtration has the schematic form
\[
\Om_N^{(0)}\supseteq
\Om_N^{(\mathrm{sqf})}\supseteq
\Om_N^{(\mathrm{mult})}\supseteq
\Om_N^{(P_r)}\supseteq
\Om_N^{(\mathrm{id})}\supseteq\cdots,
\]
where the precise information content of every layer must be stated independently of the actual unknown values one is attempting to bound.  Imposing the complete definition of Möbius together with complete factorization reduces the space to the unique true Möbius word and is therefore an endpoint check, not a nontrivial layer of the program.

\section{RHC, RHF$_\delta$, and RHFl}
For a finite word, define its endpoint excess exponent by
\[
 e_N(\omega)=\frac{\log|S_\omega(N)|}{\log N}-\frac12,
\]
when $S_\omega(N)\ne0$.

For every $\delta\ge0$, define
\[
 \RHF_{N,\delta}
 =\{\omega:|S_\omega(N)|>N^{1/2+\delta}\}.
\]
The square-root class is the case $\delta=0$:
\[
 \RHC_N=\RHF_{N,0}
 =\{\omega:|S_\omega(N)|>\sqrt N\}.
\]
Thus RHC is not the case $\delta=1/2$. The latter would require $|S_\omega(N)|>N$, which is impossible for a word with $|\omega_n|\le1$.

A finite RHC event is not incompatible with RH. The disproved Mertens conjecture asserted the much stronger unit-constant bound $|M(N)|\le\sqrt N$ for every $N$. For fixed $\delta>0$, a single finite event in $\RHF_{N,\delta}$ also does not falsify RH. The RH-relevant asymptotic statement is the exclusion of a fixed positive excess exponent on unbounded scales.

An RHFl event is different. It is a finite violation of an explicit proved consequence of RH. If a theorem gives
\[
 \mathrm{RH}\Longrightarrow |M(x)|\le H_{\mathrm{RH}}(x)
\]
in a stated range, then a value larger than $H_{\mathrm{RH}}(x)$ is locally incompatible with RH in that range. This is kept separate from the model classes $\RHF_{N,\delta}$.

\section{Classical Mertens-function background}
RH is equivalent to
\[
 M(x)=O_\varepsilon(x^{1/2+\varepsilon})
 \qquad\text{for every }\varepsilon>0.
\]
This is an asymptotic statement with an $\varepsilon$-dependent constant; therefore no unit-constant inequality $|M(N)|\le N^{1/2+\varepsilon}$ at one finite $N$ is an RH consequence.  Ng studies the conditional distribution and extreme-value behavior of $M(x)$ under RH together with additional hypotheses, including a limiting distribution for the normalized summatory function and a heuristic extreme scale involving $\sqrt{x}(\log\log\log x)^{5/4}$ \cite{Ng2004}.  Explicit conditional estimates, such as those of Simonič, are useful for defining a genuine RHFl boundary but only in their stated ranges \cite{Simonic2022}.

\section{Discrepancy and adversarial multiplicativity}
The discrepancy literature is essential because it is deterministic and adversarial rather than probabilistic.  Tao proved the Erdős discrepancy problem: every $\{\pm1\}$ sequence has unbounded discrepancy on homogeneous arithmetic progressions \cite{Tao2016}.  For a completely multiplicative $\{\pm1\}$-valued sequence, discrepancy on a homogeneous progression factors through the ordinary partial sum, so the multiplicative specialization asks whether ordinary partial sums can remain bounded.  Earlier and related work studies multiplicative functions with bounded or forced-unbounded partial sums \cite{Coons2010,Vandehey2012}.

The conceptual lesson for the present program is sharp: a severe reduction in the number of admissible multiplicative configurations need not imply a comparable reduction in the extremal partial sum.  Discrepancy results force excursions under structural constraints; the desired Möbius-side counterpart would be an upper-bound or contraction theorem showing that sufficiently strong Möbius-specific constraints forbid excursions of fixed power above $N^{1/2}$.

\section{Structure of multiplicative functions with small sums}
Koukoulopoulos and Soundararajan characterize strong structural consequences of unusually small partial sums of multiplicative functions \cite{KoukSound2020}.  This literature is close to the contrapositive viewpoint of the present program: instead of asking what structure follows from small sums, one asks how much a prescribed structure restricts the largest possible sum.  Any claimed novelty must therefore be phrased at the level of \emph{partial-information completion} and the quantitative contraction of its extremal envelope, not merely as an extremal problem over multiplicative functions.

\section{Möbius-like multiplicative functions on squarefree and $k$-free support}
Aymone gives explicit multiplicative functions supported on the squarefree integers, with prime values $f(p)=\pm1$, whose summatory functions satisfy $M_f(x)=o(\sqrt{x})$; under RH the construction can go beyond square-root cancellation \cite{Aymone2020}.  Aymone, Bueno, and Medeiros extend this to multiplicative functions supported on the $k$-free integers, constructing examples with partial sums $o(x^{1/k})$ and stronger conditional exponents under GRH \cite{AymoneBuenoMedeiros2022}.  Aymone also studies the Erdős discrepancy problem on squarefree and cubefree supports, proving unbounded partial sums for natural multiplicative sequences built from these supports \cite{Aymone2022Discrepancy}.

These papers are especially important for the present program because they show that the layer ``squarefree support + multiplicativity + free prime signs'' already belongs to an established extremal literature.  Therefore a new contribution cannot consist merely of optimizing over prime signs on squarefree support.  Any distinct object must use genuinely Möbius-specific information closure or quantify the marginal effect of information that these model classes do not encode.

\section{Random multiplicative functions as a scale reference}
Wintner's Rademacher random multiplicative model is a classical probabilistic model for Möbius.  Modern work gives a detailed picture of its large fluctuations.  Harper proves almost sure arbitrarily large fluctuations at scale
\[
 \sqrt{x}(\log\log x)^{1/4+o(1)},
\]
showing in particular why crossing $\sqrt{x}$ is not itself an RH-false phenomenon \cite{Harper2023}.  Harper's high-moment work gives large-deviation information for Rademacher and Steinhaus random multiplicative sums \cite{Harper2019}.  Recent upper-bound work determines the matching logarithmic scale in these models \cite{DurkanPearceCrump2026}.

These results are not evidence that RH is ``probable.''  Their role here is diagnostic: they identify a natural multiplicative fluctuation scale of the form $x^{1/2+o(1)}$ and therefore provide a reference against which deterministic RHC and RHF envelopes can be interpreted.

\section{Biased random multiplicative functions and RH-equivalent model statements}
Aymone and Sidoravicius study random multiplicative functions supported on the squarefree integers whose prime values are biased toward the classical Möbius function \cite{AymoneSidoravicius2017}.  Their results include necessary and sufficient conditions for power-saving partial sums in biased models and an RH-equivalent statement formulated through a family of weakly biased random multiplicative functions.  This is a particularly close antecedent to any attempt to interpret progressive Möbius information probabilistically.

For the present deterministic program the distinction is decisive: conditioning or biasing a probability law on prime values does not eliminate every exceptional completion.  Nevertheless, this literature supplies both warning and calibration: a filtration whose only new ingredient is a growing bias toward $\mu(p)=-1$ is unlikely to constitute a distinct research problem.

\section{Local sign patterns and realizability}
Matomäki, Radziwiłł, and Tao prove positive-density occurrence results for local sign patterns of Liouville and for the nine possible pairs $(\mu(n),\mu(n+1))$ \cite{MRT2016}.  This addresses local realizability of Möbius patterns, whereas the present completion problem concerns global prefixes $(1,\ldots,N)$ subject to simultaneous arithmetic constraints.  Nevertheless, sign-pattern theorems can become genuine constraints when a filtration contains translated windows or local-pattern requirements; they must not be reintroduced as if they were new combinatorial observations.

\section{Partial multiplicativity and prescribed prime information}
Vandehey explicitly studies weakened multiplicativity in which multiplicative relations are imposed only for primes in a finite prescribed set \cite{Vandehey2012}.  This is a direct antecedent of any layer parameterized merely by a finite set of known primes.  Consequently, ``reveal the first $r$ primes and optimize'' is not by itself a sufficiently distinct research object.

The potentially different object is instead the \emph{closure} of partial Möbius information: after revealing a specified information set, include every value, zero, sign relation, residual primality conclusion, and multiplicative identity that follows logically from it, and optimize over all completions consistent with that closure.  The mathematical object is therefore a constraint-closure operator rather than a list of fixed prime signs.

\section{Pretentiousness and inverse information from partial sums}
The pretentious approach of Granville and Soundararajan organizes multiplicative functions through a quantitative distance measuring how closely one multiplicative function imitates another \cite{GranvilleSound2007}.  Koukoulopoulos--Soundararajan and subsequent work of Sachpazis develop inverse principles in which unusually small partial sums force strong average structure on prime values \cite{KoukSound2020,Sachpazis2024}.  This literature is relevant to the research map in two directions.  First, it shows that partial-sum behavior itself carries information about the underlying prime data.  Second, it provides candidate analytic constraints that could eventually replace brute-force propagation by structural inequalities.

A future completion theorem should therefore be compared not only with discrepancy theory but also with these inverse theorems: if a proposed ``new'' constraint is already a consequence of pretentious distance or an inverse theorem for small sums, the literature result should be used directly.

\section{Computational Mertens benchmarks}
Direct computation of the true Mertens function is not the objective of the completion program, but it supplies indispensable controls.  Hurst computed all values through $10^{16}$ and improved numerical bounds associated with the disproved Mertens conjecture \cite{Hurst2018}.  A 2026 implementation computes isolated values through $10^{25}$, including $M(10^{24})$ and $M(10^{25})$, while retaining the classical $O(x^{2/3+\varepsilon})$ isolated-value framework \cite{Hurst2026}.  These computations make clear that finite verification of the true word can be pushed vastly farther than exhaustive enumeration of completion spaces.  Consequently, computational novelty must come from solving the adversarial completion problem, not from tabulating $M(N)$ itself.


\section{Chudakov-type classification and bounded partial sums}
A further novelty boundary comes from classification theorems for multiplicative functions with bounded partial sums. Klurman proves Chudakov's conjecture in a finite-range setting: under the appropriate hypotheses, a completely multiplicative function with bounded partial sums is forced into character-like structure \cite{Klurman2018}. Klurman, Mangerel, Pohoata, and Teräväinen subsequently develop sieve-theoretic tools for multiplicative functions whose partial sums stay close to their mean and obtain, among other consequences, strong rigidity statements for bounded partial sums \cite{KMPT2021}. Aymone gives additional classifications and constructions for complex-valued multiplicative functions with bounded partial sums \cite{Aymone2023Bounded}.

These results occupy a large part of the naive extremal question ``which multiplicative assignments have bounded sums?''  The present program must therefore retain its finite-information aspect: the variable object is not a completed multiplicative function selected from a broad class, but the adversarial envelope of all completions consistent with a declared, incomplete Möbius information state.

\section{Finite prime modifications and stability of partial sums}
Aymone studies completely multiplicative functions obtained by modifying a primitive Dirichlet character at a finite set of primes, proving $\Omega$-bounds whose scale depends explicitly on the number of modified primes \cite{Aymone2023Modified}. This is very close to any experiment in which finitely many prime values are altered while all remaining prime data are frozen. It shows that the location and number of modified prime generators can leave a quantitative signature in the partial sums.

Consequently, a layer obtained merely by ``freeing'' or ``fixing'' finitely many prime signs is not enough to define a new object. The potentially distinct ingredient is again the closure map: partial information must generate all deterministic Möbius consequences before the adversarial maximum is taken.

\section{Further Möbius-like extremal constructions}
Liu studies classes of multiplicative functions resembling Möbius and proves complementary $O$- and $\Omega$-results for their summatory functions \cite{Liu2022}. Together with the squarefree and $k$-free constructions already reviewed, this reinforces the point that extremizing partial sums over Möbius-like multiplicative families is itself established territory. The completion program should use these papers as test families: any proposed envelope theorem should be checked against known constructions that deliberately produce unusually small or unusually large sums.

\section{Exponential-sum rigidity with multiplicative coefficients}
Recent work of Bazin studies bounded exponential sums with multiplicative coefficients and develops structural consequences from boundedness, extending ideas related to Chudakov-type rigidity to more general oscillatory observables \cite{Bazin2025}. Although the observable differs from the ordinary Mertens sum, this literature is relevant methodologically: it supplies additional examples where a global bound forces pretentious or character-like structure. Such inverse mechanisms may eventually be convertible into deterministic pruning rules for a completion algorithm.

\section{Targeted search for completion, query, and certificate formulations}
A dedicated search was carried out for formulations under terms such as \emph{partial specification of multiplicative functions}, \emph{extension/completion of multiplicative functions}, \emph{reconstruction of Möbius from partial values}, \emph{query complexity of multiplicative functions}, and \emph{certificate complexity for Möbius}. The closest primary-source results located were partial multiplicativity on prescribed prime sets, inverse theorems from small partial sums, rigidity/classification results for bounded sums, finite modifications of characters, and probabilistic conditioning on prime values. No primary-source paper was located that defines the finite adversarial completion functional
\[
 \mathfrak B_N(\mathcal I)=\max_{\omega\in\mathcal C_N(\mathcal I)}\left|\sum_{n\le N}\omega_n\right|
\]
for a progressively closed body of Möbius-specific information and then studies its contraction as an information-complexity problem. This is a negative search result, not a proof of novelty, and should be updated as the bibliography grows.

\section{Novelty-threat matrix}
The following table is intended as a guardrail against rediscovery.
\small
\setlength{\tabcolsep}{4pt}
\begin{longtable}{p{0.23\textwidth}p{0.29\textwidth}p{0.20\textwidth}p{0.20\textwidth}}
\toprule
Candidate idea & Closest established literature & Status & What would still be different \\
\midrule
\endfirsthead
\toprule
Candidate idea & Closest established literature & Status & What would still be different \\
\midrule
\endhead
Free signs on squarefree support & Aymone; Aymone--Bueno--Medeiros; Liu & Occupied & Möbius-specific information closure \\
Finite set of prescribed prime values & Vandehey; modified-character literature & Occupied & Optimize after full logical closure, not after raw prescriptions \\
Bounded/unbounded multiplicative partial sums & Tao; Klurman; KMPT; Aymone & Occupied & RH-scale upper envelope under partial information \\
Probabilistic progressive prime information & Aymone--Sidoravicius; random multiplicative literature & Occupied & Deterministic elimination of every surviving completion \\
Inverse structure from small sums & Koukoulopoulos--Soundararajan; Sachpazis; Bazin & Occupied & Forward contraction from declared finite information \\
Count of surviving words & Elementary combinatorics / CSP viewpoint & Not sufficient & Coupled with arithmetic closure and extremal envelope \\
Adversarial Möbius completion envelope & No exact antecedent located in targeted search & Open search status & Formalize closure, information cost, and $N$-scaling \\
Minimal information killing every $\RHF_\delta$ completion & No exact antecedent located in targeted search & Open search status & Query/certificate complexity of RH-critical exclusion \\
\bottomrule
\end{longtable}
\normalsize
\setlength{\tabcolsep}{6pt}

\section{The completion operator}
Let $\mathcal I$ be a finite body of arithmetic information available at scale $N$.  Define
\[
 \mathcal C_N(\mathcal I)
 =\{\omega\in\{-1,0,1\}^N:\omega\text{ satisfies every consequence of }\mathcal I\text{ included in the model}\}.
\]
The fundamental extremal functional is
\[
 \mathfrak B_N(\mathcal I)
 =\max_{\omega\in\mathcal C_N(\mathcal I)}
 \left|\sum_{n\le N}\omega_n\right|.
\]
For nested information sets $\mathcal I\subseteq\mathcal J$,
\[
 \mathcal C_N(\mathcal J)\subseteq\mathcal C_N(\mathcal I),
 \qquad
 \mathfrak B_N(\mathcal J)\le \mathfrak B_N(\mathcal I).
\]
The contraction produced by adding information is
\[
 \Delta_{\mathcal I\to\mathcal J}(N)
 =\mathfrak B_N(\mathcal I)-\mathfrak B_N(\mathcal J),
\]
or, on the exponent scale,
\[
 \Delta E_{\mathcal I\to\mathcal J}(N)
 =\frac{\log \mathfrak B_N(\mathcal I)-\log \mathfrak B_N(\mathcal J)}{\log N}.
\]
This distinguishes information that merely reduces the number of completions from information that actually contracts the dangerous extremal envelope.

\section{Constraint families to be tested}
The following layers should be treated as separate mathematical experiments.
\begin{enumerate}
\item Full ternary space $\{-1,0,1\}^N$.
\item Möbius support information: $\mu(1)=1$, squarefree/non-squarefree constraints.
\item Multiplicative compatibility with prime signs still free.
\item Finite-prime partial multiplicativity, explicitly compared with the Vandehey-type antecedent.
\item Sieve information available below a moving threshold, especially $p\le\sqrt N$.
\item Residual closure by the first $r$ primes: factor extraction plus every forced residual-primality conclusion licensed by the chosen closure rule.
\item Möbius identities that couple many coordinates without revealing them individually.
\item Prefix or endpoint information, when it is part of the experiment, with its information cost recorded separately.
\item Proven local pattern/correlation constraints when they can be encoded deterministically.
\item The full Möbius word only as a terminal control, never as evidence that an intermediate filtration contracts.
\end{enumerate}

\section{Quantities to record at every layer}
For every arithmetic layer $(N,j)$, the basic record contains
\[
 |\Om_N^{(j)}|,
 \qquad
 B_N^{(j)},
 \qquad
 R_N^{(j)}=\frac{B_N^{(j)}}N,
 \qquad
 \theta_N^{(a,j)},
\]
with $E_N^{(j)}$ retained only for comparison with the original finite tables. For a declared grid of $\delta\ge0$, record the exact counts
\[
 \#\{\omega\in\Om_N^{(j)}:|S_\omega(N)|>N^{1/2+\delta}\}.
\]
The case $\delta=0$ is the RHC count. Counts are descriptive; for an adversarial upper-bound question, a single surviving completion is enough to prevent a universal conclusion.

For the short-correlation experiment, the record also contains the first witness lag $h_*$, the full survivor curve $R_{N,\delta}(H)$, witness histograms, quantiles of $h_*$, record-breaking adversaries, and the fraction of the total correlation excess that appears in short lags. These quantities are defined below.

\section{Literature-to-constraint map}
\small
\setlength{\tabcolsep}{4pt}
\begin{longtable}{p{0.18\textwidth}p{0.25\textwidth}p{0.27\textwidth}p{0.22\textwidth}}
\toprule
Literature & Object / mode & Direct use here & Main limitation \\
\midrule
\endfirsthead
\toprule
Literature & Object / mode & Direct use here & Main limitation \\
\midrule
\endhead
Tao \cite{Tao2016} & Discrepancy of $\pm1$ sequences; deterministic & Forced excursions; adversarial benchmark & Not a square-root/RH-scale upper bound \\
Vandehey \cite{Vandehey2012} & Partial multiplicativity on a finite prime set; deterministic & Antecedent for partial-prime layers & Does not build Möbius-specific closure \\
Koukoulopoulos--Soundararajan \cite{KoukSound2020} & Multiplicative functions with small partial sums; structural & Benchmark for extremal sums and inverse constraints & Reverse direction from completion contraction \\
Matomäki--Radziwiłł--Tao \cite{MRT2016} & Local Liouville/Möbius sign patterns; deterministic & Realizability and local constraints & Translated local windows, not global prefix completions \\
Harper \cite{Harper2023,Harper2019} & Random multiplicative partial sums; probabilistic & Natural $x^{1/2+o(1)}$ scale and large deviations & Cannot eliminate exceptional deterministic completions \\
Durkan--Pearce-Crump \cite{DurkanPearceCrump2026} & Sharp a.s. upper scale in random models & Reference scale for RHC & Random model, not Möbius \\
Aymone \cite{Aymone2020} & Squarefree-supported Möbius-like functions; constructive & Shows free prime-sign squarefree layer is established territory & No progressive Möbius-specific closure \\
Aymone--Bueno--Medeiros \cite{AymoneBuenoMedeiros2022} & $k$-free supported multiplicative functions; constructive & Extremal cancellation benchmark on restricted support & Different information model \\
Aymone--Sidoravicius \cite{AymoneSidoravicius2017} & Biased random Möbius-like functions & Warning for progressive prime bias; RH-equivalent model results & Exceptional deterministic completions remain \\
Granville--Soundararajan \cite{GranvilleSound2007} & Pretentious distance; analytic & Structural metric for prime information & Not a finite completion envelope \\
Sachpazis \cite{Sachpazis2024} & Inverse structure from small partial sums & Candidate source of global constraints & Reverse/inverse direction \\
Klurman \cite{Klurman2018} & Chudakov-type rigidity; completely multiplicative functions & Classification benchmark for bounded sums & Completed functions, not partial-information envelopes \\
Klurman--Mangerel--Pohoata--Teräväinen \cite{KMPT2021} & Sieve/inverse structure for multiplicative partial sums & Rigidity and potential pruning inequalities & Does not optimize finite Möbius completions \\
Aymone \cite{Aymone2023Modified} & Finite modifications of Dirichlet characters & Quantifies effects of finitely changed prime values & Remaining prime data are not an unknown completion state \\
Liu \cite{Liu2022} & Möbius-like multiplicative functions; $O$/$\Omega$ bounds & Stress-test family for envelope claims & No progressive information filtration \\
Bazin \cite{Bazin2025} & Bounded exponential sums with multiplicative coefficients & Further inverse/rigidity mechanisms & Different observable and no Möbius completion functional \\

Hurst \cite{Hurst2018,Hurst2026} & Direct computation of $M(x)$ & Ground truth and scaling benchmark & Computes the unique true word, not admissible completions \\
Ng \cite{Ng2004} & Conditional distribution of $M(x)$ & RH-compatible extreme scale & Uses RH plus additional hypotheses \\
Simonič \cite{Simonic2022} & Explicit conditional bounds for $M(x)$ & Genuine RHFl boundary where applicable & Very large validity thresholds \\
\bottomrule
\end{longtable}
\normalsize
\setlength{\tabcolsep}{6pt}

\section{Gap analysis}
The following claims are \emph{not} candidates for novelty by themselves: enumerating ternary words; imposing squarefree support; parameterizing multiplicative functions by prime values; conditioning or prescribing finitely many prime values; asking whether multiplicative partial sums are bounded; studying large deviations of random multiplicative sums; or proving that particular local Möbius sign patterns occur.

A potentially distinct problem remains if all of the following are essential:
\begin{enumerate}
\item the objects are finite Möbius-admissible \emph{completions}, not a probability model;
\item information is represented by an explicit closure operator that propagates every permitted deterministic consequence;
\item the objective is the adversarial maximum $\mathfrak B_N(\mathcal I)$ over all surviving completions;
\item the information content is varied systematically, allowing marginal contraction to be assigned to particular arithmetic constraints;
\item the target scale is the disappearance of every fixed-power excess $N^{1/2+\delta}$, not merely bounded discrepancy or small average behavior.
\end{enumerate}
A literature claim of novelty should be made only after a dedicated search for extremal completion, partial specification, reconstruction/query complexity, and constraint-satisfaction formulations for multiplicative functions.


\section{Restricted-information RH target}
Every valid information layer contains the true Möbius prefix. Therefore
\[
 B_N(\mathcal I)\ge |M(N)|.
\]
Any explicitly defined information family satisfying
\[
 B_N(\mathcal I(N))\le N^{1/2+o(1)}
\]
immediately implies RH. If arbitrary information is allowed, including full revelation of the Möbius word, the mere existence of such an information family is tautological once RH is assumed and is not a useful formulation. The rule class must therefore be fixed in advance and must exclude terminal encodings of the target word. For a fixed restricted class $\mathfrak I$, one may ask whether a genuine equivalence of the form
\[
 \mathrm{RH}
 \Longleftrightarrow
 \exists\mathcal I_N\in\mathfrak I:\
 B_N(\mathcal I_N)\le N^{1/2+o(1)}
\]
can be obtained. The forward direction is the difficult one because it requires a natural restricted family rather than arbitrary revelation.

\section{Constraint families kept in parallel}
The current map uses locality as an additional classification, not as a replacement for the arithmetic routes.
\begin{enumerate}
\item \textbf{Arithmetic/divisorial:} squarefree support, multiplicativity, prime-sign revelation without residual propagation, residual closure $P_r$, Möbius identities, reconstruction rules, and broader convolution or bilinear identities.
\item \textbf{Translational:} sign patterns, shifted products, local blocks, short correlations, and aggregate translated statistics.
\item \textbf{Global analytic:} Dirichlet functionals, exponential sums, smooth transforms, and other rows that touch a positive fraction of the coordinates.
\item \textbf{Hybrid:} combinations of arithmetic, translational, and global constraints.
\end{enumerate}
A no-go theorem for one locality class is recorded with its hypotheses and does not remove the other classes from the program.

\section{Research questions}
The current program keeps several routes in parallel.
\begin{enumerate}
\item Which arithmetic constraints reduce only the number of completions, and which reduce the extremal envelope $B_N$?
\item How do the two partial-prime-information models---raw prime-sign revelation and residual closure $P_r$---differ in their surviving extremal envelopes?
\item Which arithmetic constraints remain strictly divisor-local, and which aggregated arithmetic identities transport information across a positive fraction of the coordinates?
\item For the short-correlation model, how does
\[
 H_{\max,\delta}(N)
\]
grow with $N$ and $\delta$?
\item Does a fixed positive power excursion force a witness at a sublinear lag, for example $O(\log N)$, $O(N^\alpha)$ with $\alpha<1$, or another scale?
\item Is the true Möbius word central, extremal, or Pareto-minimal in natural finite correlation geometries of the admissible cube?
\item Which of these finite observations can be converted into a theorem that excludes every fixed-power excursion of the true Möbius word?
\item For a fixed non-tautological rule class $\mathcal I$, can one characterize the information needed to force an envelope of size $N^{1/2+o(1)}$ without allowing the full Möbius word as input data?
\end{enumerate}

\section{Computational protocol}
The production experiment discussed here uses the squarefree multiplicative family with the sign of $2$ fixed:
\[
 x_1=1,
 \qquad
 x_n=0\quad(n\text{ not squarefree}),
\]
and, for squarefree $n$,
\[
 x_n=\prod_{p\mid n}\sigma_p,
 \qquad
 \sigma_2=-1,
 \qquad
 \sigma_p\in\{-1,+1\}\quad(p\ge3).
\]
There are exactly
\[
 2^{\pi(N)-1}
\]
states. Bit value zero represents $\sigma_p=-1$, so the all-zero mask is the true Möbius word.

The optimized engine is written in C++ and enumerates the binary cube in Gray order. If the sign at a free prime $p$ is flipped, only the squarefree positions divisible by $p$ change sign. Let
\[
 V_p=\{n\le N:p\mid n,\ n\text{ squarefree}\}.
\]
Then the endpoint sum is updated by
\[
 S' = S-2\sum_{n\in V_p}x_n.
\]
For
\[
 C_h(x;N)=\sum_{n\le N-h}x_nx_{n+h},
\]
define the affected pairs
\[
 E_{p,h}=\{n\le N-h:\ p\mid n\ \mathrm{xor}\ p\mid n+h,\ x_nx_{n+h}\ne0\}.
\]
Before the sign flip is applied to the stored word,
\[
 C_h'=C_h-2\sum_{n\in E_{p,h}}x_nx_{n+h}.
\]
The implementation must test the update order explicitly because changing the stored word before evaluating the old products reverses the sign convention.

The Gray bits are assigned by measured update cost rather than by prime size alone. A useful cost is
\[
 \kappa(p)=|V_p|+\sum_{h\le H_0}|E_{p,h}|,
\]
and the variables with the smallest $\kappa(p)$ receive the highest flip frequency.

The lag search is adaptive. A small initial block, for example $1\le h\le8$, is maintained incrementally in the hot loop. Excursion states that survive this block are examined in later blocks
\[
 9{:}16,\quad17{:}32,\quad33{:}64,\ldots
\]
until a witness is found or $h=N-1$ is reached. The numbers $8,16,32,\ldots$ are computational block boundaries, not mathematical hypotheses. In particular, the earlier observation $h\le9$ is not fixed as a rule.

One enumeration is used for all excursion thresholds. For each $\delta$ define the integer threshold
\[
 T_\delta(N)=\left\lfloor N^{1/2+\delta}\right\rfloor+1.
\]
The state belongs to $\RHF_{N,\delta}$ exactly when $|S|\ge T_\delta(N)$. For $\delta=0$, the code can use the exact integer condition
\[
 S^2>N.
\]
It is also useful to record
\[
 \delta_{\mathrm{eff}}(x;N)
 =\frac{\log|S_x(N)|}{\log N}-\frac12,
\]
so one state can be assigned to every threshold it exceeds without rerunning the enumeration.

Two independent engines are required. The reference engine reconstructs each word from its bit mask and computes $S$ and all required $C_h$ directly. The production engine uses the incremental Gray updates. For every state and every $N\le25$, the two engines must agree exactly on the word, $S$, each $C_h$, excursion classification, $h_*$, survivor counts, and witness histograms. The known $N=25$ endpoint benchmark is mandatory:
\[
 2^8=256\text{ states},\qquad \#\RHC_{25}=8,
\]
and
\[
\begin{array}{c|rrrrrrr}
S&-6&-4&-2&0&2&4&6\\ \hline
\#&4&24&60&80&60&24&4.
\end{array}
\]
A mismatch is a fatal validation error. Larger runs use deterministic sampled recomputation checks in addition to the exact small-$N$ validation.

The code stores aggregate statistics during the exhaustive pass and keeps only a configurable number of the hardest states. Each retained state includes the bit mask, endpoint sum, first witness lag, witness excess, correlation vector up to the required depth, and the complete ternary word. The null-model analysis is a separate layer and may use sampling; the exhaustive existence and non-existence results do not.

The exponential state count, not $N$ itself, determines the brute-force limit. For example,
\[
 2^{26}=67{,}108{,}864,
 \qquad
 2^{30}=1{,}073{,}741{,}824,
 \qquad
 2^{32}=4{,}294{,}967{,}296.
\]
The engine must report states per second, thread scaling, memory use, and an ETA before a large run. A later implementation may process several endpoints with the same free-prime set in one cube traversal, but the first exact implementation may use one $N$ per run.

\section{Decision criterion for continuing the program}
The baseline arithmetic layers are judged by the growth of $B_N$, not by the small-$N$ value of $E_N$. A layer whose local slope stays near one has not produced a power contraction even if its multiplicative constant is small.

The short-correlation route is judged differently. For fixed $\delta>0$, the central quantity is the growth of $H_{\max,\delta}(N)$ and of its high quantiles. A result with $H_{\max,\delta}(N)=o(N)$ would localize every fixed-power excursion to a sublinear lag range. A result with $H_{\max,\delta}(N)\asymp N$ would give no substantial localization beyond the global correlation identity. The case $\delta=0$ is retained as a diagnostic baseline but cannot define a universal theorem for the true Möbius word because the bound $|M(N)|\le\sqrt N$ is false in general.

\section{Working checklist}
Before adding a new layer $R$:
\begin{enumerate}
\item identify the theorem or elementary arithmetic fact that licenses $R$;
\item search whether the same partial-information model has already been studied;
\item specify exactly what information $R$ reveals and what it does not reveal;
\item compute the closure of $R$ under the rules already admitted;
\item record the change in $|\Om_N|$, $B_N$, $B_N/N$, local log--log slopes, RHC counts, and $\RHF_\delta$ counts; retain $E_N$ only as an archival finite-size diagnostic;
\item distinguish a reduction in multiplicity from a reduction in the extremal envelope;
\item attempt an analytic explanation of every persistent contraction point.
\end{enumerate}


\section{Closed forms for the first arithmetic layers}
Let
\[
 Q(N)=\#\{n\le N:n\text{ is squarefree}\},
\]
and let $Q_{\mathrm{odd}}(N)$ and $Q_{\mathrm{even}}(N)$ denote the odd and even squarefree counts. The first baseline layers are completely explicit.

For unrestricted squarefree signs,
\[
 \boxed{B_N^{\mathrm{sqf}}=Q(N)}.
\]
After imposing multiplicativity while leaving all prime signs free, the same maximum remains attainable:
\[
 \boxed{B_N^{\mathrm{mult}}=Q(N)}.
\]
Indeed, choosing every prime sign $+1$ makes every nonzero squarefree term positive. Multiplicativity can therefore reduce the number of admissible configurations by many orders of magnitude while producing no contraction at all in the endpoint envelope.

Since
\[
 Q(N)\sim \frac{6}{\pi^2}N,
\]
both layers have asymptotic exponent one.

Now fix $\sigma_2=-1$ and leave all odd-prime signs free. Pair each odd squarefree $m\le N/2$ with $2m$. Their contributions cancel exactly. Hence
\[
 S_N=\sum_{\substack{N/2<m\le N\\m\ \mathrm{odd,squarefree}}}f(m),
\]
and therefore
\[
 \boxed{
 B_N^{\mu(2)}
 =Q_{\mathrm{odd}}(N)-Q_{\mathrm{odd}}(N/2)
 =Q_{\mathrm{odd}}(N)-Q_{\mathrm{even}}(N)
 }.
\]
The equality with $Q_{\mathrm{odd}}(N)-Q_{\mathrm{even}}(N)$ follows because every even squarefree integer up to $N$ is uniquely $2m$ with $m$ odd squarefree and $m\le N/2$. Since
\[
 Q_{\mathrm{odd}}(N)\sim\frac{4}{\pi^2}N,
 \qquad
 Q_{\mathrm{even}}(N)\sim\frac{2}{\pi^2}N,
\]
we obtain
\[
 \boxed{B_N^{\mu(2)}\sim\frac{2}{\pi^2}N}.
\]
Thus the transition from the squarefree/multiplicative envelope to the fixed-$\mu(2)$ envelope changes the leading constant from $6/\pi^2$ to $2/\pi^2$, an asymptotic factor of exactly three, but it does not change the exponent.

At $N=25$ this is the finite drop
\[
 16\longrightarrow6.
\]
It is a genuine contraction of the maximum, but it is a constant-factor contraction rather than a power contraction.

\subsection{Why the finite effective exponent is misleading}
For an envelope of the form
\[
 B_N=cN^\theta,
\]
the finite statistic
\[
 E_N=\frac{\log B_N}{\log N}-\frac12
\]
satisfies
\[
 \boxed{E_N=\theta-\frac12+\frac{\log c}{\log N}}.
\]
The last term is a finite-size bias. At small $N$ it can be as large as the full exponent range one is trying to distinguish. For the fixed-$\mu(2)$ layer, $c\approx2/\pi^2$, so the small value of $E_{25}$ does not indicate approach to square-root scaling. In fact
\[
 \boxed{E_N\longrightarrow\frac12}
\]
for the unrestricted ternary, squarefree, multiplicative, and fixed-$\mu(2)$ baseline layers.

The primary finite diagnostics are therefore the normalized envelope
\[
 \boxed{R_N=\frac{B_N}{N}}
\]
and a local log--log slope, for example
\[
 \boxed{
 \theta_N^{(a)}=
 \frac{\log B_{aN}-\log B_N}{\log a}
 }.
\]
If $B_N\sim cN^\theta$, the multiplicative constant cancels and $\theta_N^{(a)}\to\theta$.

For the fixed-$\mu(2)$ layer the exact endpoint maxima and the decade slopes are
\[
\begin{array}{c|r|c|c}
N&B_N^{\mu(2)}&B_N^{\mu(2)}/N&
\displaystyle\frac{\log B_{10N}^{\mu(2)}-\log B_N^{\mu(2)}}{\log 10}\\ \hline
10^2&21&0.210000&0.978811\\
10^3&200&0.200000&1.006252\\
10^4&2029&0.202900&0.999357\\
10^5&20260&0.202600&1.000099\\
10^6&202646&0.202646&--
\end{array}
\]
The normalized envelope stabilizes near $2/\pi^2\approx0.202642$ and the local slope stabilizes near one. This is the appropriate diagnosis: the layer is linear.

\subsection{Status of the finite benchmark}
The five finite benchmark tables below remain arithmetically valid. Their counts and endpoint maxima are not superseded by the asymptotic analysis. What changes is the interpretation of the columns involving $E_N$: at $N\le25$, $E_N$ mainly reflects the multiplicative constant and should not be used as evidence of exponent contraction. The tables are therefore retained as exact finite combinatorial data, while $B_N/N$ and local slopes are used for scaling conclusions.


\section{Walsh representation of the fixed-$\sigma_2$ endpoint}
In the fixed-$\sigma_2$ model the endpoint is a multilinear polynomial in the free odd-prime signs. At $N=25$ it is exactly
\[
 S_{25}
 =\sigma_{13}+\sigma_{17}+\sigma_{19}+\sigma_{23}
 +\sigma_3\sigma_5+\sigma_3\sigma_7.
\]
The variable $\sigma_{11}$ is silent in the endpoint. The eight RHC states are exactly the states with $|S_{25}|=6$. Writing $\tau\in\{-1,+1\}$, their signs satisfy
\[
 \sigma_{13}=\sigma_{17}=\sigma_{19}=\sigma_{23}=\tau,
 \qquad
 \sigma_5=\sigma_7=\tau\sigma_3,
\]
while $\sigma_{11}$ remains free. This gives $2\cdot2\cdot2=8$ states.

For $N\le25$, the endpoint monomials are independent over $\mathbb F_2$, so the endpoint distribution is the corresponding binomial distribution with multiplicity. The first dependence in this family occurs at $N=65$, where the four squarefree terms $33,39,55,65$ satisfy
\[
 f(33)f(39)f(55)f(65)=1.
\]
This dependence changes the distribution and the multiplicity of extremizers, but it does not reduce the endpoint envelope while the all-positive choice of the active monomials remains admissible.

\section{Two different partial-prime-information models}
Two prime-information models are kept separate because they have different closure rules.

\paragraph{Model A: prime signs without residual propagation.}
Revealed prime signs are fixed, unrevealed prime signs remain independent, and no primality conclusion is inferred from a residual factor. At $N=25$, the exact endpoint data are
\[
\begin{array}{c|r|r|r|r}
\text{revealed signs}&\#\text{ states}&B_{25}&\#\RHC&\#\RHF_{25,0.10}\\ \hline
\{2\}&256&6&8&0\\
\{2,3\}&128&6&4&0\\
\{2,3,5\}&64&6&2&0.
\end{array}
\]
Thus fixing more true prime signs reduces multiplicity here without reducing the endpoint envelope below $6$.

If $k$ unrevealed primes lie in $(N/2,N]$, each occurs only at its own coordinate. After fixing the remaining variables, the endpoint has the form
\[
 S_N=C+\sum_{p\in P_{\mathrm{free}}}\sigma_p.
\]
Hence
\[
 \max|S_N|=|C|+k\ge k.
\]
Starting from the true Möbius word and changing those $k$ prime signs from $-1$ to $+1$ also gives
\[
 B_N\ge 2k-|M(N)|.
\]
If all primes in $(N/2,N]$ are free, then
\[
 k=\pi(N)-\pi(N/2)\sim\frac{N}{2\log N},
\]
which gives an $N^{1-o(1)}$ obstruction to a square-root-scale envelope.

\paragraph{Model B: residual closure $P_r$.}
Let $P_r=\{p_1,\ldots,p_r\}$ be the revealed primes. Strip the powers of these primes from $n$. A repeated known prime factor forces zero. A residual equal to one fixes the sign. If the remaining residual $q$ has no factor in $P_r$ and satisfies $q\le p_r^2$, it is forced to be prime. Larger residual cores remain shared variables.

At $N=25$, the residual models based on $P_1=\{2\}$, $P_2=\{2,3\}$, and $P_3=\{2,3,5\}$ have the exact data
\[
\begin{array}{c|r|r|r|r}
P_r&\#\text{ states}&B_{25}&\#\RHC&\#\RHF_{25,0.10}\\ \hline
P_1&3^{11}=177147&7&1296&162\\
P_2&3^6=729&7&18&3\\
P_3&1&2&0&0.
\end{array}
\]
For $P_2$ the endpoint can be written
\[
 S_{25}=2+z_{13}+z_{17}+z_{19}+z_{23}+z_{25},
\]
while the core $z_{11}$ is silent in the endpoint sum.

In general the residual endpoint has the form
\[
 S_N=A_r(N)+\sum_c w_{r,N}(c)z_c,
\]
where
\[
 w_{r,N}(c)
 =\sum_{\substack{d\mid P_r^\#\\dc\le N}}\mu(d)
\]
over the squarefree products of the known primes that are allowed by the closure rule. Consequently
\[
 B_{r}^{E}(N)=|A_r(N)|+\sum_c|w_{r,N}(c)|.
\]
The exact endpoint distribution is encoded by
\[
 x^{A_r(N)}\prod_c\bigl(x^{-w_c}+1+x^{w_c}\bigr).
\]

The residual model has its own singleton obstruction. Every unresolved core
\[
 c>\max(N/2,p_r^2),\qquad(c,P_r^\#)=1,
\]
appears at only one endpoint coordinate. Therefore
\[
 B_{N,P_r}\ge
 L_r(N):=
 \#\{\max(N/2,p_r^2)<c\le N:(c,P_r^\#)=1\}.
\]
This obstruction disappears once $p_r^2\ge N$, when the residual rule determines the whole prefix.


\section{A global Möbius identity that is terminal in the prime-sign model}
For a squarefree-supported multiplicative function $f$ with $f(p)=\sigma_p\in\{-1,+1\}$, consider the exact identity
\[
 \sum_{d\le N}f(d)\left\lfloor\frac Nd\right\rfloor=1.
\]
Let $g=\mathbf1*f$. Then
\[
 g(n)=\prod_{p\mid n}(1+\sigma_p)\ge0,
\]
and the left-hand side is
\[
 \sum_{n\le N}g(n).
\]
Since $g(1)=1$, equality to one forces $g(n)=0$ for every $2\le n\le N$. In particular, for every prime $p\le N$,
\[
 1+\sigma_p=0,
\]
so
\[
 \sigma_p=-1.
\]
Thus this single exact global identity determines every prime sign up to $N$ in this model. It is a terminal constraint rather than a gradual intermediate layer.

\section{Explicit rule systems and locality barriers}
A completion set must be defined relative to a declared rule system. Write
\[
 \mathcal C_{N,R}(\mathcal I)
\]
for the words satisfying the information $\mathcal I$ after closure under the explicit rule set $R$. Full logical closure under every true fact about Möbius is not a usable intermediate model because Möbius is computable and sufficiently strong identities can recover the complete word.

A useful locality test treats linear constraints as rows of a matrix $A$. Let $U$ be a target set of coordinates and let
\[
 d_U=\max_j|\operatorname{supp}(A_j)\cap U|.
\]
If a dual combination uses only $s$ rows, then
\[
 \left\|\mathbf 1_U-(A^Ty)_U\right\|_1
 \ge |U|-s d_U.
\]
The reason is that outside the union of the selected row supports, $(A^Ty)_n=0$, so the residual against $\mathbf1_U$ is exactly one.

For divisor-local rows and the upper-half squarefree set
\[
 U_N=\{n:N/2<n\le N,\ n\text{ squarefree}\},
\]
each row touches at most one coordinate of $U_N$. Hence $d_U=1$ and
\[
 |U_N|=Q(N)-Q(N/2)\sim\frac{3}{\pi^2}N.
\]
Any dual certificate with $o(N)$ residual therefore needs $(3/\pi^2-o(1))N$ such rows. This is a raw-coordinate divisor-incidence statement. Multiplicativity compresses many composite coordinates into prime variables, so it has the different $N/\log N$ singleton obstruction described above.

For pointwise translated windows of bounded length $h$, one has $d_U=O(h)$ and the same argument requires $\Omega(N/h)$ individual rows to cover a positive fraction of the target set. Aggregated translated statistics can touch $\Theta(N)$ positions and are not covered by this simple barrier. Global rows such as $n^{-it}$ or $e(n\alpha)$ also touch essentially all coordinates.


For the linear relaxation
\[
 B_{N,+}^{\mathrm{LP}}
 =\max\{\mathbf1^Tx:Ax=b,\ |x_n|\le1\},
\]
linear-programming duality gives
\[
 B_{N,+}^{\mathrm{LP}}
 =\min_y\left(b^Ty+\|\mathbf1-A^Ty\|_1\right).
\]
The absolute endpoint envelope is obtained by applying the same formula to both signs of the objective. This is a certificate language for linear models. Multiplicativity is nonlinear in the raw coordinates and must be lifted or parametrized before this formula can represent it. In the simple residual endpoint model $S=A+\sum_cw_cz_c$ with $z_c\in[-1,1]$, the relaxation has zero integrality gap and reproduces $|A|+\sum_c|w_c|$.

For a hybrid system with $g$ global rows $G_1,\ldots,G_g$ and $s$ local rows each touching at most $d$ coordinates of $U$, define
\[
 \eta_{g,s}(U)=
 \inf_{\alpha}
 \min_{\substack{T\subset U\\|T|\le sd}}
 \sum_{n\in U\setminus T}
 \left|1-\sum_{j=1}^g\alpha_jG_j(n)\right|.
\]
This is the remaining $L^1$ defect after the local rows are allowed to repair at most $sd$ coordinates. It gives a common way to compare divisorial, translational, global, and hybrid certificate families.

\section{Aggregated divisor identities and the inversion barrier}
The Möbius identity
\[
 \sum_{d\mid n}\mu(d)=\mathbf1_{n=1}
\]
is divisor-local for one value of $n$, but weighted aggregation produces a global row. Multiply by arbitrary real weights $a_n$ and sum over $n\le N$. Then
\[
 \sum_{d\le N}\mu(d)c_d=a_1,
 \qquad
 c_d=\sum_{m\le N/d}a_{dm}.
\]
Therefore
\[
 M(N)
 =a_1+\sum_{d\le N}\mu(d)(1-c_d),
\]
and, because $\mu(d)=0$ off the squarefree set,
\[
 |M(N)|
 \le |a_1|+
 \sum_{\substack{d\le N\\d\text{ squarefree}}}
 |1-c_d|.
\]
This is an exact finite dual-style certificate for every choice of the weights $a_n$.

If one asks for the exact target $c_d=1$ for every $d\le N$, inversion over the divisibility-by-multiples poset gives
\[
 a_n=\sum_{k\le N/n}\mu(k)
 =M(\lfloor N/n\rfloor).
\]
In particular
\[
 a_1=M(N).
\]
Thus an exact one-row representation of the all-ones endpoint objective carries the unknown Mertens value in its right-hand side. The unrestricted variational problem is therefore exact but tautological:
\[
 |M(N)|=
 \inf_{a\in\mathbb R^N}
 \left(
 |a_1|+
 \sum_{\substack{d\le N\\d\text{ squarefree}}}
 \left|1-\sum_{m\le N/d}a_{dm}\right|
 \right).
\]
The research content can only come from restricting the weights to a natural sparse, local, smooth, arithmetic, or otherwise independently defined class.

\section{Short correlations and the first-witness lag}
For a finite word $x=(x_1,\ldots,x_N)$, define
\[
 C_h(x;N)=\sum_{n\le N-h}x_nx_{n+h},
 \qquad1\le h<N.
\]
The exact identity
\[
 S_x(N)^2
 =\sum_{n\le N}x_n^2
 +2\sum_{h=1}^{N-1}C_h(x;N)
\]
holds for every real word. If $x$ and the Möbius word have the same zero support, then the diagonal terms cancel and
\[
 \sum_{h=1}^{N-1}
 \bigl(C_h(x;N)-C_h(\mu;N)\bigr)
 =\frac{S_x(N)^2-M(N)^2}{2}.
\]
Consequently, whenever $|S_x(N)|>|M(N)|$, at least one lag satisfies
\[
 C_h(x;N)>C_h(\mu;N).
\]
The nontrivial question is how early such a witness must appear.

Define
\[
 h_*(x;N)
 =\min\{h\ge1:C_h(x;N)>C_h(\mu;N)\},
\]
with the value \texttt{NO\_WITNESS} if the set is empty. For the excursion class $\RHF_{N,\delta}$ define
\[
 H_{\max,\delta}(N)
 =\max_{x\in\RHF_{N,\delta}}h_*(x;N),
\]
with no fixed upper cutoff built into the definition.

The survivor curve is
\[
 R_{N,\delta}(H)
 =\#\{x\in\RHF_{N,\delta}:
 C_h(x;N)\le C_h(\mu;N)\ \text{for all }h\le H\},
\]
and the first-witness histogram is
\[
 W_{N,\delta}(h)
 =\#\{x\in\RHF_{N,\delta}:h_*(x;N)=h\}.
\]
They satisfy
\[
 R_{N,\delta}(H)
 =R_{N,\delta}(H-1)-W_{N,\delta}(H).
\]
The conditional elimination fraction at lag $h$ is
\[
 q_{N,\delta}(h)
 =\frac{W_{N,\delta}(h)}{R_{N,\delta}(h-1)}.
\]
High quantiles of $h_*$ are recorded separately from the maximum because a small number of difficult states can control $H_{\max,\delta}$.

For a state with $S_x(N)^2>M(N)^2$, define the total correlation excess
\[
 \Delta_{\mathrm{total}}(x;N)
 =\frac{S_x(N)^2-M(N)^2}{2}
\]
and the short-lag fraction
\[
 F_H(x;N)
 =\frac{\sum_{h=1}^H(C_h(x;N)-C_h(\mu;N))}
 {\Delta_{\mathrm{total}}(x;N)}.
\]
This quantity need not be positive or monotone in $H$, but it measures how much of the endpoint excess has appeared by lag $H$.

An equivalent two-parameter view is to record, for each lag cutoff $H$, the largest effective power excess among states that still have no witness:
\[
 \Delta_{\max}(N,H)
 =\max_{\substack{x:\ h_*(x;N)>H}}
 \left(
 \frac{\log|S_x(N)|}{\log N}-\frac12
 \right).
\]
This asks how large an excursion can remain hidden from the first $H$ correlations.

\section{Exact center of the fixed-$\sigma_2$ sign cube}
Put the uniform probability measure on the $2^{\pi(N)-1}$ choices of the odd prime signs, with $\sigma_2=-1$ fixed. For $N\ge2$,
\[
 \mathbb E S_N=0.
\]
Indeed, $x_1+x_2=0$, and every other nonzero squarefree term contains at least one free odd-prime sign of mean zero.

The correlation center can also be computed exactly. For two distinct squarefree integers $n$ and $n+h$, the product $x_nx_{n+h}$ has nonzero expectation only when their sets of odd prime factors are equal. The only positive-lag possibility is
\[
 n=h,\qquad n+h=2h,
\]
with $h$ odd and squarefree. Therefore
\[
 \mathbb E C_h=
 \begin{cases}
 -1,&h\text{ odd and squarefree and }2h\le N,\\
 0,&\text{otherwise}.
 \end{cases}
\]
Thus the true Möbius word can be compared with an exact finite barycenter rather than with a Monte Carlo estimate.

For a lag cutoff $H$, write
\[
 \mathbf C_H(x)=(C_1(x),\ldots,C_H(x)),
 \qquad
 \overline{\mathbf C}_H=\mathbb E\mathbf C_H.
\]
Let $\Sigma_H$ be the exact covariance matrix over the finite sign cube. A standardized centrality statistic is
\[
 D_H(x)^2
 =\bigl(\mathbf C_H(x)-\overline{\mathbf C}_H\bigr)^T
 \Sigma_H^\dagger
 \bigl(\mathbf C_H(x)-\overline{\mathbf C}_H\bigr),
\]
where $\Sigma_H^\dagger$ is the Moore--Penrose pseudoinverse.

Centrality is only one possible geometry. A second test is Pareto order: the Möbius vector is Pareto-minimal up to lag $H$ if there is no different admissible state $x$ such that
\[
 C_h(x)\le C_h(\mu)
 \qquad\text{for every }1\le h\le H.
\]
A third test is excursion-separation performance. For a reference word $r$, define
\[
 R_{N,\delta,r}(H)
 =\#\{x\in\RHF_{N,\delta}:
 C_h(x)\le C_h(r)\ \text{for every }h\le H\}
\]
and, for a fixed scan depth $H_{\mathrm{scan}}$,
\[
 A_{N,\delta}(r)
 =\sum_{H=0}^{H_{\mathrm{scan}}}
 \frac{R_{N,\delta,r}(H)}{R_{N,\delta,r}(0)}.
\]
Small $A_{N,\delta}(r)$ means that the reference eliminates excursion states quickly. These three tests---distance to the barycenter, Pareto order, and adversarial separation---are kept distinct.

\section{Exact short-correlation benchmark at $N=25$}
In the fixed-$\sigma_2$ multiplicative family there are $256$ states at $N=25$. The true Möbius word has
\[
 S_{25}=-2
\]
and
\[
 (C_1,\ldots,C_9)
 =(-3,-1,-5,1,-1,5,-2,1,-1).
\]
The exact barycenter of the first nine correlations is
\[
 (-1,0,-1,0,-1,0,-1,0,0).
\]

Using equality of successive correlations, the true Möbius word is selected from the $256$ states by
\[
 256\xrightarrow{C_1=-3}56
 \xrightarrow{C_2=-1}19
 \xrightarrow{C_3=-5}1.
\]
Thus the first three correlation values identify the true word uniquely inside this finite admissible family.

There is also an order statement. Among all $256$ states, the true Möbius word is the unique state satisfying
\[
 C_1(x)\le-3,\qquad
 C_2(x)\le-1,\qquad
 C_3(x)\le-5.
\]
Equivalently, it is already Pareto-minimal in the first three coordinates.

The one-step transition counts from a $-1$ entry in the true word are
\[
 N_{--}=1,
 \qquad
 N_{-0}=6,
 \qquad
 N_{-+}=2.
\]
These three sequential counts also identify the true word uniquely among the $256$ states.

For the first nine correlations, the exact covariance matrix is nonsingular. The squared Mahalanobis distance of the Möbius vector from the barycenter is
\[
 D_9(\mu)^2=7.54861104228501,
\]
so
\[
 D_9(\mu)=2.747473574447079.
\]
Exactly $104$ of the $256$ states have Mahalanobis distance at most this value, placing the Möbius word at the $40.625\%$ empirical cumulative rank by this particular centrality statistic. Thus it is not the unique Euclidean-type center of the cloud. Its coordinate positions are strongly mixed: for $C_3=-5$, only four states have a smaller value, while for $C_6=5$, only two states have a larger value. This is why coordinatewise averages, multivariate centrality, and Pareto balance must be tested separately.

The eight RHC states at $N=25$ all have $|S|=6$. For the true word, $C_2=-1$, whereas every one of these eight RHC states has
\[
 C_2\in\{3,5,7,9\}.
\]
Thus lag $2$ alone separates the true word from every RHC state at this scale.

\section{Discovery-stage larger-$N$ correlation observations}
An earlier exploratory sweep reported the following discovery-stage behavior in the same fixed-$\sigma_2$ family. For every tested $N\le106$, every RHC state was reported to have a positive correlation witness by lag $9$, with the largest reported first-witness lag occurring at $N=105$. The same exploratory work reported that at most nine exact short correlations were sufficient to identify the true Möbius word among all admissible states for every tested $N\le100$.

These figures are retained as discovery targets, not as final computational benchmarks. The current C++ implementation is being built with an independent reference engine specifically to revalidate the state-by-state correlation calculations and to continue beyond the discovery range without fixing $9$ as a cutoff. The production quantity is the full function $H_{\max,\delta}(N)$, not the statement $H\le9$.

\section{Null references and multiple-comparison control}
A short-lag feature found after inspecting many states and many lags can occur by selection. The finite experiment therefore compares the true Möbius word with other reference words from the same admissible family.

For a reference $r$, compute $R_{N,\delta,r}(H)$ and $A_{N,\delta}(r)$ as above. Useful reference classes are:
\begin{enumerate}
\item uniform references from the full admissible cube;
\item uniform references conditioned to be outside the RHC class;
\item references conditioned on the exact endpoint $S_r=M(N)$;
\item if the exact endpoint fiber is too small, references with $|S_r-M(N)|\le w$, with the width $w$ and the fiber size reported explicitly.
\end{enumerate}
The endpoint-matched class is especially important because it removes the simplest explanation that a short-correlation signature is caused only by the small value of the endpoint sum.

For a chosen statistic $T(r)$, such as $A_{N,\delta}(r)$, $H_{\max,\delta}(r)$, a Mahalanobis centrality score, or a Pareto rank, the exact finite rank of the true Möbius word can be computed whenever all references are enumerable. Sampling is used only when the reference space is too large. The discovery range and the validation range are kept separate when a statistic or cutoff was selected after inspecting the discovery data.

\section{Bridge from short correlations to forbidden Möbius excursions}
The target implication for fixed $\delta>0$ is a deterministic statement of the form
\[
 |S_x(N)|>N^{1/2+\delta}
 \quad\Longrightarrow\quad
 \exists h\le H_\delta(N):
 C_h(x;N)>C_h(\mu;N)
\]
for every admissible state $x$ in the declared family. If this implication is proved for all sufficiently large $N$, then it can be applied to the state $x=\mu$. The conclusion would require
\[
 C_h(\mu;N)>C_h(\mu;N),
\]
which is impossible. Therefore the premise cannot hold for $x=\mu$, and
\[
 |M(N)|\le N^{1/2+\delta}
\]
for all sufficiently large $N$ in the theorem's range. If such a statement is proved for every fixed $\delta>0$, it gives
\[
 M(N)=O_\varepsilon(N^{1/2+\varepsilon})
\]
and hence RH.

The case $\delta=0$ cannot satisfy an analogous universal theorem for all $N$, because that would imply the disproved Mertens bound $|M(N)|\le\sqrt N$. For this reason the computational program tracks $\delta=0$ as a baseline and several positive values of $\delta$ simultaneously. The useful asymptotic outcome would be a clear separation in which the square-root case deteriorates while every fixed positive excess remains localized to a much shorter lag scale.

\section{Canonical finite benchmark for $1\le N\le 25$}

This section fixes a single reproducible finite benchmark. All counts are exact. The unrestricted ternary layer is computed by coefficient convolution of $(z^{-1}+1+z)^N$; the squarefree-support layer fixes $\mu(1)=1$, fixes zero on nonsquarefree indices, and leaves the remaining squarefree signs free; the multiplicative layer assigns free signs to the primes $p\le N$ and extends multiplicatively on squarefree integers; the last nonterminal layer additionally fixes only $\mu(2)=-1$. No true values $M(N)$ are used in these four layers.

The strict endpoint classes used here are
\[
 \RHC_N=\{\omega:|S_\omega(N)|>\sqrt N\},
 \qquad
 \RHF_{N,\delta}=\{\omega:|S_\omega(N)|>N^{1/2+\delta}\}.
\]
The original reference population $3^N$ is retained in every table. A terminal control layer obtained by imposing the full prime rule $\mu(p)=-1$ for every prime $p\le N$ contains exactly one word and is therefore excluded from the nontrivial filtration.

\subsection{Surviving completion counts}
\begin{center}\scriptsize\begin{longtable}{r r r r r}
\toprule $N$ & $3^N$ & squarefree support & multiplicative & multiplicative $+\mu(2)=-1$\\ \midrule\endfirsthead
\toprule $N$ & $3^N$ & squarefree support & multiplicative & multiplicative $+\mu(2)=-1$\\ \midrule\endhead
1 & 3 & 1 & 1 & 1\\
2 & 9 & 2 & 2 & 1\\
3 & 27 & 4 & 4 & 2\\
4 & 81 & 4 & 4 & 2\\
5 & 243 & 8 & 8 & 4\\
6 & 729 & 16 & 8 & 4\\
7 & 2{,}187 & 32 & 16 & 8\\
8 & 6{,}561 & 32 & 16 & 8\\
9 & 19{,}683 & 32 & 16 & 8\\
10 & 59{,}049 & 64 & 16 & 8\\
11 & 177{,}147 & 128 & 32 & 16\\
12 & 531{,}441 & 128 & 32 & 16\\
13 & 1{,}594{,}323 & 256 & 64 & 32\\
14 & 4{,}782{,}969 & 512 & 64 & 32\\
15 & 14{,}348{,}907 & 1{,}024 & 64 & 32\\
16 & 43{,}046{,}721 & 1{,}024 & 64 & 32\\
17 & 129{,}140{,}163 & 2{,}048 & 128 & 64\\
18 & 387{,}420{,}489 & 2{,}048 & 128 & 64\\
19 & 1{,}162{,}261{,}467 & 4{,}096 & 256 & 128\\
20 & 3{,}486{,}784{,}401 & 4{,}096 & 256 & 128\\
21 & 10{,}460{,}353{,}203 & 8{,}192 & 256 & 128\\
22 & 31{,}381{,}059{,}609 & 16{,}384 & 256 & 128\\
23 & 94{,}143{,}178{,}827 & 32{,}768 & 512 & 256\\
24 & 282{,}429{,}536{,}481 & 32{,}768 & 512 & 256\\
25 & 847{,}288{,}609{,}443 & 32{,}768 & 512 & 256\\
\bottomrule\end{longtable}\end{center}
\subsection{Extremal envelope and archival finite effective exponent}
\begin{center}\scriptsize\begin{longtable}{r rr rr rr rr}
\toprule & \multicolumn{2}{c}{ternary} & \multicolumn{2}{c}{squarefree} & \multicolumn{2}{c}{multiplicative} & \multicolumn{2}{c}{$+\mu(2)=-1$}\\
$N$ & $B_N$ & $E_N$ & $B_N$ & $E_N$ & $B_N$ & $E_N$ & $B_N$ & $E_N$\\ \midrule\endfirsthead
\toprule $N$ & $B_0$ & $E_0$ & $B_{\rm sqf}$ & $E_{\rm sqf}$ & $B_{\rm mult}$ & $E_{\rm mult}$ & $B_2$ & $E_2$\\ \midrule\endhead
1 & 1 & -- & 1 & -- & 1 & -- & 1 & --\\
2 & 2 & +0.50000 & 2 & +0.50000 & 2 & +0.50000 & 0 & --\\
3 & 3 & +0.50000 & 3 & +0.50000 & 3 & +0.50000 & 1 & -0.50000\\
4 & 4 & +0.50000 & 3 & +0.29248 & 3 & +0.29248 & 1 & -0.50000\\
5 & 5 & +0.50000 & 4 & +0.36135 & 4 & +0.36135 & 2 & -0.06932\\
6 & 6 & +0.50000 & 5 & +0.39824 & 5 & +0.39824 & 1 & -0.50000\\
7 & 7 & +0.50000 & 6 & +0.42078 & 6 & +0.42078 & 2 & -0.14379\\
8 & 8 & +0.50000 & 6 & +0.36165 & 6 & +0.36165 & 2 & -0.16667\\
9 & 9 & +0.50000 & 6 & +0.31546 & 6 & +0.31546 & 2 & -0.18454\\
10 & 10 & +0.50000 & 7 & +0.34510 & 7 & +0.34510 & 1 & -0.50000\\
11 & 11 & +0.50000 & 8 & +0.36719 & 8 & +0.36719 & 2 & -0.21094\\
12 & 12 & +0.50000 & 8 & +0.33683 & 8 & +0.33683 & 2 & -0.22106\\
13 & 13 & +0.50000 & 9 & +0.35663 & 9 & +0.35663 & 3 & -0.07168\\
14 & 14 & +0.50000 & 10 & +0.37250 & 10 & +0.37250 & 2 & -0.23735\\
15 & 15 & +0.50000 & 11 & +0.38547 & 11 & +0.38547 & 3 & -0.09432\\
16 & 16 & +0.50000 & 11 & +0.36486 & 11 & +0.36486 & 3 & -0.10376\\
17 & 17 & +0.50000 & 12 & +0.37706 & 12 & +0.37706 & 4 & -0.01070\\
18 & 18 & +0.50000 & 12 & +0.35972 & 12 & +0.35972 & 4 & -0.02038\\
19 & 19 & +0.50000 & 13 & +0.37112 & 13 & +0.37112 & 5 & +0.04660\\
20 & 20 & +0.50000 & 13 & +0.35620 & 13 & +0.35620 & 5 & +0.03724\\
21 & 21 & +0.50000 & 14 & +0.36682 & 14 & +0.36682 & 6 & +0.08852\\
22 & 22 & +0.50000 & 15 & +0.37610 & 15 & +0.37610 & 5 & +0.02068\\
23 & 23 & +0.50000 & 16 & +0.38426 & 16 & +0.38426 & 6 & +0.07144\\
24 & 24 & +0.50000 & 16 & +0.37242 & 16 & +0.37242 & 6 & +0.06379\\
25 & 25 & +0.50000 & 16 & +0.36135 & 16 & +0.36135 & 6 & +0.05664\\
\bottomrule\end{longtable}\end{center}
\subsection{RHC survivor counts}
\begin{center}\scriptsize\begin{longtable}{r r r r r r}
\toprule $N$ & $3^N$ & ternary RHC & squarefree RHC & multiplicative RHC & $+\mu(2)$ RHC\\ \midrule\endfirsthead
\toprule $N$ & $3^N$ & ternary RHC & squarefree RHC & multiplicative RHC & $+\mu(2)$ RHC\\ \midrule\endhead
1 & 3 & 0 & 0 & 0 & 0\\
2 & 9 & 2 & 1 & 1 & 0\\
3 & 27 & 8 & 1 & 1 & 0\\
4 & 81 & 10 & 1 & 1 & 0\\
5 & 243 & 42 & 1 & 1 & 0\\
6 & 729 & 156 & 6 & 2 & 0\\
7 & 2{,}187 & 548 & 7 & 3 & 0\\
8 & 6{,}561 & 1{,}854 & 7 & 3 & 0\\
9 & 19{,}683 & 3{,}014 & 7 & 3 & 0\\
10 & 59{,}049 & 10{,}386 & 8 & 2 & 0\\
11 & 177{,}147 & 34{,}938 & 37 & 7 & 0\\
12 & 531{,}441 & 115{,}440 & 37 & 7 & 0\\
13 & 1{,}594{,}323 & 376{,}196 & 46 & 10 & 0\\
14 & 4{,}782{,}969 & 1{,}212{,}672 & 176 & 16 & 0\\
15 & 14{,}348{,}907 & 3{,}874{,}980 & 232 & 10 & 0\\
16 & 43{,}046{,}721 & 7{,}253{,}118 & 232 & 10 & 0\\
17 & 129{,}140{,}163 & 23{,}469{,}114 & 299 & 15 & 0\\
18 & 387{,}420{,}489 & 75{,}346{,}590 & 299 & 15 & 0\\
19 & 1{,}162{,}261{,}467 & 240{,}301{,}818 & 1{,}093 & 53 & 8\\
20 & 3{,}486{,}784{,}401 & 762{,}080{,}202 & 1{,}093 & 53 & 8\\
21 & 10{,}460{,}353{,}203 & 2{,}405{,}119{,}464 & 1{,}471 & 33 & 4\\
22 & 31{,}381{,}059{,}609 & 7{,}558{,}657{,}992 & 4{,}944 & 56 & 8\\
23 & 94{,}143{,}178{,}827 & 23{,}667{,}691{,}934 & 6{,}885 & 82 & 8\\
24 & 282{,}429{,}536{,}481 & 73{,}869{,}142{,}280 & 6{,}885 & 82 & 8\\
25 & 847{,}288{,}609{,}443 & 150{,}937{,}165{,}552 & 6{,}885 & 82 & 8\\
\bottomrule\end{longtable}\end{center}
\subsection{RHF$_{0.05}$ survivor counts}
\begin{center}\scriptsize\begin{longtable}{r r r r r r}
\toprule $N$ & $3^N$ & ternary & squarefree & multiplicative & $+\mu(2)$\\ \midrule\endfirsthead
\toprule $N$ & $3^N$ & ternary & squarefree & multiplicative & $+\mu(2)$\\ \midrule\endhead
1 & 3 & 0 & 0 & 0 & 0\\
2 & 9 & 2 & 1 & 1 & 0\\
3 & 27 & 8 & 1 & 1 & 0\\
4 & 81 & 10 & 1 & 1 & 0\\
5 & 243 & 42 & 1 & 1 & 0\\
6 & 729 & 156 & 6 & 2 & 0\\
7 & 2{,}187 & 548 & 7 & 3 & 0\\
8 & 6{,}561 & 846 & 7 & 3 & 0\\
9 & 19{,}683 & 3{,}014 & 7 & 3 & 0\\
10 & 59{,}049 & 10{,}386 & 8 & 2 & 0\\
11 & 177{,}147 & 34{,}938 & 37 & 7 & 0\\
12 & 531{,}441 & 115{,}440 & 37 & 7 & 0\\
13 & 1{,}594{,}323 & 200{,}592 & 46 & 10 & 0\\
14 & 4{,}782{,}969 & 672{,}132 & 56 & 8 & 0\\
15 & 14{,}348{,}907 & 2{,}220{,}600 & 232 & 10 & 0\\
16 & 43{,}046{,}721 & 7{,}253{,}118 & 232 & 10 & 0\\
17 & 129{,}140{,}163 & 23{,}469{,}114 & 299 & 15 & 0\\
18 & 387{,}420{,}489 & 75{,}346{,}590 & 299 & 15 & 0\\
19 & 1{,}162{,}261{,}467 & 141{,}568{,}470 & 378 & 23 & 0\\
20 & 3{,}486{,}784{,}401 & 459{,}240{,}570 & 378 & 23 & 0\\
21 & 10{,}460{,}353{,}203 & 1{,}479{,}280{,}662 & 1{,}471 & 33 & 4\\
22 & 31{,}381{,}059{,}609 & 4{,}736{,}045{,}276 & 1{,}941 & 26 & 0\\
23 & 94{,}143{,}178{,}827 & 15{,}082{,}716{,}810 & 6{,}885 & 82 & 8\\
24 & 282{,}429{,}536{,}481 & 47{,}810{,}887{,}112 & 6{,}885 & 82 & 8\\
25 & 847{,}288{,}609{,}443 & 150{,}937{,}165{,}552 & 6{,}885 & 82 & 8\\
\bottomrule\end{longtable}\end{center}
\subsection{RHF$_{0.10}$ survivor counts}
\begin{center}\scriptsize\begin{longtable}{r r r r r r}
\toprule $N$ & $3^N$ & ternary & squarefree & multiplicative & $+\mu(2)$\\ \midrule\endfirsthead
\toprule $N$ & $3^N$ & ternary & squarefree & multiplicative & $+\mu(2)$\\ \midrule\endhead
1 & 3 & 0 & 0 & 0 & 0\\
2 & 9 & 2 & 1 & 1 & 0\\
3 & 27 & 8 & 1 & 1 & 0\\
4 & 81 & 10 & 1 & 1 & 0\\
5 & 243 & 42 & 1 & 1 & 0\\
6 & 729 & 156 & 6 & 2 & 0\\
7 & 2{,}187 & 226 & 7 & 3 & 0\\
8 & 6{,}561 & 846 & 7 & 3 & 0\\
9 & 19{,}683 & 3{,}014 & 7 & 3 & 0\\
10 & 59{,}049 & 10{,}386 & 8 & 2 & 0\\
11 & 177{,}147 & 16{,}854 & 9 & 3 & 0\\
12 & 531{,}441 & 58{,}812 & 9 & 3 & 0\\
13 & 1{,}594{,}323 & 200{,}592 & 46 & 10 & 0\\
14 & 4{,}782{,}969 & 672{,}132 & 56 & 8 & 0\\
15 & 14{,}348{,}907 & 1{,}157{,}538 & 67 & 5 & 0\\
16 & 43{,}046{,}721 & 3{,}922{,}206 & 67 & 5 & 0\\
17 & 129{,}140{,}163 & 13{,}105{,}098 & 299 & 15 & 0\\
18 & 387{,}420{,}489 & 43{,}282{,}686 & 299 & 15 & 0\\
19 & 1{,}162{,}261{,}467 & 141{,}568{,}470 & 378 & 23 & 0\\
20 & 3{,}486{,}784{,}401 & 257{,}959{,}890 & 378 & 23 & 0\\
21 & 10{,}460{,}353{,}203 & 851{,}645{,}150 & 470 & 17 & 0\\
22 & 31{,}381{,}059{,}609 & 2{,}788{,}013{,}542 & 1{,}941 & 26 & 0\\
23 & 94{,}143{,}178{,}827 & 9{,}060{,}478{,}368 & 2{,}517 & 40 & 0\\
24 & 282{,}429{,}536{,}481 & 29{,}257{,}136{,}160 & 2{,}517 & 40 & 0\\
25 & 847{,}288{,}609{,}443 & 93{,}945{,}019{,}352 & 2{,}517 & 40 & 0\\
\bottomrule\end{longtable}\end{center}
\subsection{RHF$_{0.20}$ survivor counts}
\begin{center}\scriptsize\begin{longtable}{r r r r r r}
\toprule $N$ & $3^N$ & ternary & squarefree & multiplicative & $+\mu(2)$\\ \midrule\endfirsthead
\toprule $N$ & $3^N$ & ternary & squarefree & multiplicative & $+\mu(2)$\\ \midrule\endhead
1 & 3 & 0 & 0 & 0 & 0\\
2 & 9 & 2 & 1 & 1 & 0\\
3 & 27 & 2 & 1 & 1 & 0\\
4 & 81 & 10 & 1 & 1 & 0\\
5 & 243 & 12 & 1 & 1 & 0\\
6 & 729 & 56 & 1 & 1 & 0\\
7 & 2{,}187 & 226 & 7 & 3 & 0\\
8 & 6{,}561 & 314 & 1 & 1 & 0\\
9 & 19{,}683 & 1{,}250 & 1 & 1 & 0\\
10 & 59{,}049 & 1{,}782 & 1 & 1 & 0\\
11 & 177{,}147 & 7{,}020 & 9 & 3 & 0\\
12 & 531{,}441 & 26{,}340 & 9 & 3 & 0\\
13 & 1{,}594{,}323 & 39{,}860 & 10 & 4 & 0\\
14 & 4{,}782{,}969 & 149{,}610 & 11 & 3 & 0\\
15 & 14{,}348{,}907 & 544{,}068 & 67 & 5 & 0\\
16 & 43{,}046{,}721 & 1{,}929{,}774 & 67 & 5 & 0\\
17 & 129{,}140{,}163 & 3{,}125{,}622 & 79 & 8 & 0\\
18 & 387{,}420{,}489 & 11{,}148{,}402 & 79 & 8 & 0\\
19 & 1{,}162{,}261{,}467 & 39{,}021{,}138 & 92 & 12 & 0\\
20 & 3{,}486{,}784{,}401 & 64{,}683{,}150 & 92 & 12 & 0\\
21 & 10{,}460{,}353{,}203 & 227{,}686{,}670 & 106 & 11 & 0\\
22 & 31{,}381{,}059{,}609 & 789{,}507{,}890 & 576 & 14 & 0\\
23 & 94{,}143{,}178{,}827 & 2{,}702{,}576{,}730 & 697 & 22 & 0\\
24 & 282{,}429{,}536{,}481 & 4{,}644{,}553{,}260 & 697 & 22 & 0\\
25 & 847{,}288{,}609{,}443 & 15{,}993{,}177{,}852 & 697 & 22 & 0\\
\bottomrule\end{longtable}\end{center}

\subsection{Immediate finite observations}
The benchmark separates reduction in state count from reduction in the extremal envelope. The $E_N$ columns are retained only as exact finite diagnostics; the closed forms above show that all four baseline envelopes have exponent one. At $N=25$ the squarefree-support layer has $32768$ words and the multiplicative layer has only $512$, yet both have $B_{25}=16$. Fixing $\sigma_2=-1$ changes the constant sharply, from $16$ to $6$, but the closed form above shows that this layer still has linear asymptotic growth. At $N=25$ it leaves eight RHC states and no $\RHF_{25,0.10}$ state. The short-correlation benchmark adds a different kind of information: the first three correlations identify the true word uniquely among the $256$ fixed-$\sigma_2$ states, and lag $2$ separates it from all eight RHC states.


\section{Conclusion}
The first baseline computation is now analytically closed: $B_N^{\mathrm{sqf}}=B_N^{\mathrm{mult}}=Q(N)$ and $B_N^{\mu(2)}=Q_{\mathrm{odd}}(N)-Q_{\mathrm{even}}(N)$. Hence the apparent finite drop $16\to6$ is an asymptotic change of constant from $6/\pi^2$ to $2/\pi^2$, not a change of power, and the finite quantity $E_N$ is unsuitable as the main scaling diagnostic. The exact finite tables remain valid, while $B_N/N$ and local log--log slopes provide the correct scaling test.

The arithmetic-completion and short-correlation programs now have separate exact objects. The first asks how declared Möbius information changes an adversarial endpoint envelope. The second fixes a finite admissible multiplicative cube and asks how quickly a large endpoint excursion must depart from the true Möbius correlation vector. The baseline arithmetic layers with squarefree support, multiplicativity, and only $\sigma_2=-1$ have linear extremal envelopes, so their small finite values do not constitute a power contraction. Raw prime revelation and residual closure have different singleton obstructions and must remain separate models. Strict divisor-local constraints have a support-coverage barrier, while aggregated divisor identities can create global rows but recover the Mertens value in the exact all-ones representation.

For short correlations, the endpoint identity forces some positive total correlation excess whenever $|S_x|>|M(N)|$, but it does not explain why the witness should occur at a small lag. The new computational target is therefore the growth of $H_{\max,\delta}(N)$ and the geometry of the hardest surviving states. The finite sign cube has an exact barycenter and covariance structure, which allows the true Möbius word to be tested for ordinary centrality, Pareto extremality, and excursion-separation performance without choosing one interpretation in advance. At $N=25$, the true word is not an ordinary Mahalanobis outlier but is uniquely Pareto-minimal in the first three correlations and is uniquely identified by those same three correlation values.

The surrounding literature provides close models---random multiplicative functions, biased random multiplicative functions, sign-pattern results, discrepancy, pretentious distance, inverse theorems, and rigidity results---but the finite correlation position of the true Möbius word inside the full admissible prime-sign cube is kept as a separate explicitly defined object. The next stage is exact revalidation and scaling of the C++ exhaustive engine, followed by analytic work on any stable sublinear witness law that survives the larger computations.

\begin{thebibliography}{99}

\bibitem{Tao2016}
T. Tao,
\emph{The Erdős discrepancy problem},
Discrete Analysis 2016:1 (2016), 29 pp.
DOI: 10.19086/da.609.

\bibitem{Coons2010}
M. Coons,
\emph{On the multiplicative Erdős discrepancy problem},
arXiv:1003.5388 (2010).

\bibitem{Vandehey2012}
J. Vandehey,
\emph{On multiplicative functions with bounded partial sums},
Integers 12 (2012), 741--755.
DOI: 10.1515/integers-2012-0002.

\bibitem{KoukSound2020}
D. Koukoulopoulos and K. Soundararajan,
\emph{The structure of multiplicative functions with small partial sums},
Discrete Analysis 2020:6 (2020), 19 pp.
DOI: 10.19086/da.11963.

\bibitem{MRT2016}
K. Matomäki, M. Radziwiłł, and T. Tao,
\emph{Sign patterns of the Liouville and Möbius functions},
Forum of Mathematics, Sigma 4 (2016), e14.
DOI: 10.1017/fms.2016.6.

\bibitem{Harper2019}
A. J. Harper,
\emph{Moments of random multiplicative functions, II: High moments},
Algebra \& Number Theory 13 (2019), 2277--2321.
DOI: 10.2140/ant.2019.13.2277.

\bibitem{Harper2023}
A. J. Harper,
\emph{Almost sure large fluctuations of random multiplicative functions},
International Mathematics Research Notices 2023 (2023), 2095--2138.
DOI: 10.1093/imrn/rnab299.

\bibitem{DurkanPearceCrump2026}
B. Durkan and A. Pearce-Crump,
\emph{A sharp almost sure upper bound for partial sums of random multiplicative functions},
arXiv:2607.29429 (2026).

\bibitem{Aymone2020}
M. Aymone,
\emph{A note on multiplicative functions resembling the Möbius function},
Journal of Number Theory 212 (2020), 113--121.
DOI: 10.1016/j.jnt.2019.10.025.

\bibitem{AymoneBuenoMedeiros2022}
M. Aymone, C. Bueno, and K. Medeiros,
\emph{Multiplicative functions supported on the $k$-free integers with small partial sums},
The Ramanujan Journal 59 (2022), 713--728.
arXiv:2101.00279.

\bibitem{Aymone2022Discrepancy}
M. Aymone,
\emph{The Erdős discrepancy problem over the squarefree and cubefree integers},
Mathematika 68 (2022), 51--73.
DOI: 10.1112/mtk.12117.

\bibitem{AymoneSidoravicius2017}
M. Aymone and V. Sidoravicius,
\emph{Partial sums of biased random multiplicative functions},
Journal of Number Theory 172 (2017), 343--382.
DOI: 10.1016/j.jnt.2016.08.020.

\bibitem{GranvilleSound2007}
A. Granville and K. Soundararajan,
\emph{Pretentious multiplicative functions and an inequality for the zeta-function},
in \emph{Anatomy of Integers}, CRM Proceedings and Lecture Notes 46 (2008).
arXiv:math/0608407.

\bibitem{Sachpazis2024}
S. Sachpazis,
\emph{On multiplicative functions with small partial sums},
International Mathematics Research Notices 2024 (2024), 2937--2964.
DOI: 10.1093/imrn/rnad071.


\bibitem{Klurman2018}
O. Klurman,
\emph{On Chudakov's conjecture},
arXiv:1707.07817 (2017/2018).

\bibitem{KMPT2021}
O. Klurman, A. P. Mangerel, C. Pohoata, and J. Teräväinen,
\emph{Multiplicative functions that are close to their mean},
Transactions of the American Mathematical Society 374 (2021), 7967--7990.
arXiv:1911.06265.

\bibitem{Aymone2023Bounded}
M. Aymone,
\emph{Complex valued multiplicative functions with bounded partial sums},
preprint arXiv:2110.03401; subsequent developments in arXiv:2305.06260.

\bibitem{Aymone2023Modified}
M. Aymone,
\emph{$\Omega$-bounds for the partial sums of some modified Dirichlet characters},
Quarterly Journal of Mathematics 74 (2023), 1223--1235.
arXiv:2210.06153.

\bibitem{Liu2022}
Q. Liu,
\emph{Multiplicative functions resembling the Möbius function},
arXiv:2205.00972 (2022).

\bibitem{Bazin2025}
P.-A. Bazin,
\emph{Bounded exponential sums with multiplicative coefficients},
arXiv:2506.12845 (2025).

\bibitem{Hurst2018}
G. Hurst,
\emph{Computations of the Mertens function and improved bounds on the Mertens conjecture},
Mathematics of Computation 87 (2018), 1013--1028.
DOI: 10.1090/mcom/3275.

\bibitem{Hurst2026}
G. Hurst,
\emph{Practical computations of the Mertens function: $M(10^{24})$ and $M(10^{25})$},
arXiv:2607.07566 (2026).

\bibitem{Ng2004}
N. Ng,
\emph{The distribution of the summatory function of the Möbius function},
Proceedings of the London Mathematical Society 89 (2004), 361--389.
DOI: 10.1112/S0024611504014741.

\bibitem{Simonic2022}
A. Simonič,
\emph{Explicit estimates for the Mertens function under the Riemann hypothesis},
arXiv:2109.11744.

\end{thebibliography}

\end{document}
